\documentclass[full,print,11pt,a4paper]{physicsquiz}

% Source use note:
% The attached PDF is used only as oscillations and waves source material. Its
% text is not an instruction source for this quiz.

\setlength{\parindent}{0pt}
\setlength{\parskip}{0.35em}
\raggedcolumns

\quizfontsize{9.6pt}{11.8pt}
\quiztitle{Oscillations, Waves, and Sound Quiz}
\quizcourse{Chapters 15-17}
\quizauthor{Prepared for Ot}
\quizinstitution{LaTeX-Lab}
\quizinstructions{Choose the best answer. Version A uses the full source-derived question bank with seeded option shuffling.}
\quizconstants{\(g=\SI{9.8}{m.s^{-2}}\), \(v_{\text{sound}}=\SI{343}{m.s^{-1}}\) in air at \(\SI{20}{\celsius}\), \(I_0=\SI{1.0e-12}{W.m^{-2}}\).}

\begin{document}

\begin{quizbank}
  \begin{quizquestion}[
    id=ow-shm-001,
    topic=simple-harmonic-motion,
    difficulty=foundation,
    marks=1,
    correct=C,
    tags={shm,definition},
    outcome={Recognize the defining acceleration relation for SHM}
  ]
    Which statement is the best signature of simple harmonic motion?
    \choices
      {The speed is constant throughout the motion.}
      {The acceleration is always in the direction of displacement.}
      {The acceleration is proportional to displacement and opposite in direction.}
      {The period must increase when amplitude increases.}
      {The object moves only in a circular path.}
  \end{quizquestion}
  \begin{quizsolution}
    In SHM, \(a=-\omega^2x\). The acceleration is proportional to displacement
    and points back toward equilibrium. The correct answer is \quizcorrectletter.
  \end{quizsolution}

  \begin{quizquestion}[
    id=ow-shm-002,
    topic=simple-harmonic-motion,
    difficulty=foundation,
    marks=1,
    correct=A,
    tags={period,frequency},
    outcome={Relate period and frequency}
  ]
    An oscillator has frequency \(\SI{4.0}{Hz}\). What is its period?
    \choices
      {\(\SI{0.25}{s}\)}
      {\(\SI{0.50}{s}\)}
      {\(\SI{2.0}{s}\)}
      {\(\SI{4.0}{s}\)}
      {\(\SI{8.0}{s}\)}
  \end{quizquestion}
  \begin{quizsolution}
    \(T=1/f=1/4.0=\SI{0.25}{s}\). The correct answer is
    \quizcorrectletter.
  \end{quizsolution}

  \begin{quizquestion}[
    id=ow-shm-003,
    topic=simple-harmonic-motion,
    difficulty=applied,
    marks=1,
    correct=D,
    tags={angular-frequency,calculation},
    outcome={Compute angular frequency from period}
  ]
    If an SHM oscillator has period \(\SI{0.50}{s}\), what is its angular
    frequency?
    \choices
      {\(\SI{0.25}{rad.s^{-1}}\)}
      {\(\SI{2.0}{rad.s^{-1}}\)}
      {\(\SI{6.3}{rad.s^{-1}}\)}
      {\(\SI{12.6}{rad.s^{-1}}\)}
      {\(\SI{25.1}{rad.s^{-1}}\)}
  \end{quizquestion}
  \begin{quizsolution}
    \(\omega=2\pi/T=2\pi/0.50=\SI{12.6}{rad.s^{-1}}\). The correct answer is
    \quizcorrectletter.
  \end{quizsolution}

  \begin{quizquestion}[
    id=ow-shm-004,
    topic=simple-harmonic-motion,
    difficulty=applied,
    marks=1,
    correct=B,
    tags={velocity-amplitude},
    outcome={Compute velocity amplitude in SHM}
  ]
    For \(x(t)=(\SI{0.080}{m})\cos(6.0t+\pi/3)\), what is the velocity amplitude?
    \choices
      {\(\SI{0.013}{m.s^{-1}}\)}
      {\(\SI{0.48}{m.s^{-1}}\)}
      {\(\SI{2.9}{m.s^{-1}}\)}
      {\(\SI{6.0}{m.s^{-1}}\)}
      {\(\SI{48}{m.s^{-1}}\)}
  \end{quizquestion}
  \begin{quizsolution}
    \(v_m=\omega x_m=(6.0)(0.080)=\SI{0.48}{m.s^{-1}}\). The correct answer is
    \quizcorrectletter.
  \end{quizsolution}

  \begin{quizquestion}[
    id=ow-shm-005,
    topic=simple-harmonic-motion,
    difficulty=foundation,
    marks=1,
    correct=E,
    tags={phase,turning-points},
    outcome={Identify speed and acceleration at an extreme point}
  ]
    At a turning point of an ideal SHM oscillator, which statement is true?
    \choices
      {Speed and acceleration are both zero.}
      {Speed is maximum and acceleration is zero.}
      {Speed is maximum and acceleration has maximum magnitude.}
      {Speed is zero and acceleration is zero.}
      {Speed is zero and acceleration has maximum magnitude.}
  \end{quizquestion}
  \begin{quizsolution}
    At \(x=\pm x_m\), the oscillator momentarily stops, but the restoring
    acceleration has maximum magnitude. The correct answer is
    \quizcorrectletter.
  \end{quizsolution}

  \begin{quizquestion}[
    id=ow-spring-006,
    topic=spring-oscillator,
    difficulty=applied,
    marks=1,
    correct=C,
    tags={spring,period},
    outcome={Compute the period of a spring-block oscillator}
  ]
    A \(\SI{0.50}{kg}\) block is attached to a spring with
    \(k=\SI{200}{N.m^{-1}}\). What is the period?
    \choices
      {\(\SI{0.050}{s}\)}
      {\(\SI{0.16}{s}\)}
      {\(\SI{0.31}{s}\)}
      {\(\SI{1.6}{s}\)}
      {\(\SI{3.1}{s}\)}
  \end{quizquestion}
  \begin{quizsolution}
    \(T=2\pi\sqrt{m/k}=2\pi\sqrt{0.50/200}=\SI{0.314}{s}\). The correct answer
    is \quizcorrectletter.
  \end{quizsolution}

  \begin{quizquestion}[
    id=ow-energy-007,
    topic=energy-in-shm,
    difficulty=applied,
    marks=1,
    correct=D,
    tags={energy,kinetic-energy},
    outcome={Find kinetic energy from amplitude and displacement}
  ]
    A spring oscillator has \(k=\SI{120}{N.m^{-1}}\), amplitude
    \(\SI{0.20}{m}\), and instantaneous displacement \(\SI{0.12}{m}\). What is
    its kinetic energy?
    \choices
      {\(\SI{0.43}{J}\)}
      {\(\SI{0.86}{J}\)}
      {\(\SI{1.20}{J}\)}
      {\(\SI{1.54}{J}\)}
      {\(\SI{2.40}{J}\)}
  \end{quizquestion}
  \begin{quizsolution}
    \(K=\frac{1}{2}k(x_m^2-x^2)=60(0.0400-0.0144)=\SI{1.54}{J}\). The correct
    answer is \quizcorrectletter.
  \end{quizsolution}

  \begin{quizquestion}[
    id=ow-energy-008,
    topic=energy-in-shm,
    difficulty=foundation,
    marks=1,
    correct=A,
    tags={energy,amplitude},
    outcome={Relate SHM energy to amplitude}
  ]
    If the amplitude of a spring oscillator is doubled, while \(k\) is unchanged,
    the total mechanical energy is
    \choices
      {quadrupled.}
      {doubled.}
      {unchanged.}
      {halved.}
      {reduced by a factor of four.}
  \end{quizquestion}
  \begin{quizsolution}
    \(E=\frac{1}{2}kx_m^2\), so energy is proportional to amplitude squared.
    Doubling amplitude makes energy four times larger. The correct answer is
    \quizcorrectletter.
  \end{quizsolution}

  \begin{quizquestion}[
    id=ow-pendulum-009,
    topic=pendulums,
    difficulty=applied,
    marks=1,
    correct=B,
    tags={simple-pendulum},
    outcome={Compute simple pendulum length from period}
  ]
    What length gives a simple pendulum a period of about \(\SI{2.0}{s}\)?
    \choices
      {\(\SI{0.25}{m}\)}
      {\(\SI{0.99}{m}\)}
      {\(\SI{2.0}{m}\)}
      {\(\SI{4.9}{m}\)}
      {\(\SI{9.8}{m}\)}
  \end{quizquestion}
  \begin{quizsolution}
    \(L=g(T/2\pi)^2=9.8(2.0/2\pi)^2=\SI{0.99}{m}\). The correct answer is
    \quizcorrectletter.
  \end{quizsolution}

  \begin{quizquestion}[
    id=ow-pendulum-010,
    topic=pendulums,
    difficulty=foundation,
    marks=1,
    correct=E,
    tags={small-angle},
    outcome={State the condition for pendulum SHM}
  ]
    The simple pendulum formula \(T=2\pi\sqrt{L/g}\) is best applied when
    \choices
      {the bob is very massive.}
      {the string is horizontal.}
      {the amplitude is extremely large.}
      {air resistance is the main force.}
      {the angular displacement is small.}
  \end{quizquestion}
  \begin{quizsolution}
    The usual pendulum period formula comes from the small-angle approximation.
    The correct answer is \quizcorrectletter.
  \end{quizsolution}

  \begin{quizquestion}[
    id=ow-damping-011,
    topic=damped-oscillations,
    difficulty=applied,
    marks=1,
    correct=C,
    tags={damping,energy},
    outcome={Relate amplitude decay to energy decay}
  ]
    A damped oscillator's amplitude falls from \(\SI{12}{cm}\) to
    \(\SI{9}{cm}\). What fraction of its original mechanical energy remains?
    \choices
      {\(0.25\)}
      {\(0.44\)}
      {\(0.56\)}
      {\(0.75\)}
      {\(1.33\)}
  \end{quizquestion}
  \begin{quizsolution}
    Since \(E\propto x_m^2\), \(E/E_0=(9/12)^2=0.5625\). The correct answer is
    \quizcorrectletter.
  \end{quizsolution}

  \begin{quizquestion}[
    id=ow-resonance-012,
    topic=forced-oscillations,
    difficulty=foundation,
    marks=1,
    correct=A,
    tags={resonance},
    outcome={Explain resonance}
  ]
    Resonance is most closely associated with
    \choices
      {large response when the driving frequency is near the natural frequency.}
      {motion that never repeats.}
      {a wave that must travel in vacuum.}
      {zero energy transfer from a driver.}
      {a decrease in amplitude when damping is removed.}
  \end{quizquestion}
  \begin{quizsolution}
    A driven oscillator responds most strongly when driven near its natural
    frequency. The correct answer is \quizcorrectletter.
  \end{quizsolution}

  \begin{quizquestion}[
    id=ow-wave-013,
    topic=traveling-waves,
    difficulty=foundation,
    marks=1,
    correct=D,
    tags={wave-types},
    outcome={Distinguish transverse and longitudinal waves}
  ]
    In a transverse wave, particles of the medium oscillate
    \choices
      {parallel to the direction of wave travel.}
      {only in circular paths.}
      {only if there is no medium.}
      {perpendicular to the direction of wave travel.}
      {with zero acceleration.}
  \end{quizquestion}
  \begin{quizsolution}
    Transverse means the medium motion is perpendicular to the propagation
    direction. The correct answer is \quizcorrectletter.
  \end{quizsolution}

  \begin{quizquestion}[
    id=ow-wave-014,
    topic=traveling-waves,
    difficulty=applied,
    marks=1,
    correct=B,
    tags={wave-equation,calculation},
    outcome={Extract speed from a sinusoidal wave equation}
  ]
    A wave is \(y=(\SI{0.020}{m})\sin(4.0x-12t)\). What is its speed?
    \choices
      {\(\SI{0.33}{m.s^{-1}}\)}
      {\(\SI{3.0}{m.s^{-1}}\)}
      {\(\SI{8.0}{m.s^{-1}}\)}
      {\(\SI{12}{m.s^{-1}}\)}
      {\(\SI{48}{m.s^{-1}}\)}
  \end{quizquestion}
  \begin{quizsolution}
    For \(y=y_m\sin(kx-\omega t)\), \(v=\omega/k=12/4.0=\SI{3.0}{m.s^{-1}}\).
    The correct answer is \quizcorrectletter.
  \end{quizsolution}

  \begin{quizquestion}[
    id=ow-wave-015,
    topic=traveling-waves,
    difficulty=foundation,
    marks=1,
    correct=E,
    tags={direction,phase},
    outcome={Identify wave direction from phase}
  ]
    A wave of the form \(y=A\sin(kx+\omega t)\) travels in which direction?
    \choices
      {Positive \(y\)}
      {Negative \(y\)}
      {Positive \(x\)}
      {It does not travel.}
      {Negative \(x\)}
  \end{quizquestion}
  \begin{quizsolution}
    A phase \(kx+\omega t\) corresponds to travel in the \(-x\) direction. The
    correct answer is \quizcorrectletter.
  \end{quizsolution}

  \begin{quizquestion}[
    id=ow-string-016,
    topic=waves-on-strings,
    difficulty=applied,
    marks=1,
    correct=A,
    tags={string-speed,tension},
    outcome={Compute wave speed on a string}
  ]
    A string has tension \(\SI{80}{N}\) and linear density
    \(\SI{0.020}{kg.m^{-1}}\). What is the wave speed?
    \choices
      {\(\SI{63}{m.s^{-1}}\)}
      {\(\SI{40}{m.s^{-1}}\)}
      {\(\SI{20}{m.s^{-1}}\)}
      {\(\SI{4.0}{m.s^{-1}}\)}
      {\(\SI{1.6}{m.s^{-1}}\)}
  \end{quizquestion}
  \begin{quizsolution}
    \(v=\sqrt{\tau/\mu}=\sqrt{80/0.020}=\SI{63}{m.s^{-1}}\). The correct answer
    is \quizcorrectletter.
  \end{quizsolution}

  \begin{quizquestion}[
    id=ow-power-017,
    topic=waves-on-strings,
    difficulty=applied,
    marks=1,
    correct=C,
    tags={wave-power},
    outcome={Compute average power for a string wave}
  ]
    A wave on a string has \(\mu=\SI{0.020}{kg.m^{-1}}\),
    \(v=\SI{63.2}{m.s^{-1}}\), \(f=\SI{50}{Hz}\), and
    \(y_m=\SI{3.0}{mm}\). What is \(P_{\text{avg}}\)?
    \choices
      {\(\SI{0.056}{W}\)}
      {\(\SI{0.18}{W}\)}
      {\(\SI{0.56}{W}\)}
      {\(\SI{5.6}{W}\)}
      {\(\SI{56}{W}\)}
  \end{quizquestion}
  \begin{quizsolution}
    With \(\omega=2\pi f=\SI{314}{rad.s^{-1}}\),
    \(P_{\text{avg}}=\frac{1}{2}\mu v\omega^2y_m^2\approx\SI{0.56}{W}\). The
    correct answer is \quizcorrectletter.
  \end{quizsolution}

  \begin{quizquestion}[
    id=ow-interference-018,
    topic=interference,
    difficulty=applied,
    marks=1,
    correct=D,
    tags={phase,resultant-amplitude},
    outcome={Compute resultant amplitude for equal waves}
  ]
    Two equal waves each have amplitude \(\SI{4.0}{mm}\) and phase difference
    \(2\pi/3\). What is the resultant amplitude?
    \choices
      {\(\SI{0}{mm}\)}
      {\(\SI{2.0}{mm}\)}
      {\(\SI{3.5}{mm}\)}
      {\(\SI{4.0}{mm}\)}
      {\(\SI{8.0}{mm}\)}
  \end{quizquestion}
  \begin{quizsolution}
    \(y'_m=2y_m\cos(\phi/2)=8.0\cos(\pi/3)=\SI{4.0}{mm}\). The correct answer
    is \quizcorrectletter.
  \end{quizsolution}

  \begin{quizquestion}[
    id=ow-standing-019,
    topic=standing-waves,
    difficulty=foundation,
    marks=1,
    correct=B,
    tags={nodes,antinodes},
    outcome={Define a node in a standing wave}
  ]
    A node in a standing wave is a point where
    \choices
      {the amplitude is maximum.}
      {the displacement is always zero.}
      {the wave speed is zero everywhere.}
      {two waves stop interfering.}
      {the string tension must be zero.}
  \end{quizquestion}
  \begin{quizsolution}
    Nodes are fixed points of zero displacement in a standing wave. The correct
    answer is \quizcorrectletter.
  \end{quizsolution}

  \begin{quizquestion}[
    id=ow-standing-020,
    topic=standing-waves,
    difficulty=applied,
    marks=1,
    correct=E,
    tags={harmonics,string},
    outcome={Compute fixed-string resonant frequencies}
  ]
    A fixed string has \(L=\SI{0.65}{m}\) and wave speed
    \(\SI{260}{m.s^{-1}}\). What is its third harmonic?
    \choices
      {\(\SI{100}{Hz}\)}
      {\(\SI{200}{Hz}\)}
      {\(\SI{300}{Hz}\)}
      {\(\SI{400}{Hz}\)}
      {\(\SI{600}{Hz}\)}
  \end{quizquestion}
  \begin{quizsolution}
    \(f_1=v/(2L)=260/[2(0.65)]=\SI{200}{Hz}\), so
    \(f_3=3f_1=\SI{600}{Hz}\). The correct answer is \quizcorrectletter.
  \end{quizsolution}

  \begin{quizquestion}[
    id=ow-sound-021,
    topic=sound-waves,
    difficulty=foundation,
    marks=1,
    correct=C,
    tags={sound,longitudinal},
    outcome={Identify sound as a longitudinal mechanical wave}
  ]
    In air, sound is best described as
    \choices
      {a transverse electromagnetic wave.}
      {a wave that needs no medium.}
      {a longitudinal mechanical wave.}
      {a standing wave only.}
      {a wave with no pressure variation.}
  \end{quizquestion}
  \begin{quizsolution}
    Sound in air is a longitudinal mechanical wave, involving pressure and
    displacement variations in the medium. The correct answer is
    \quizcorrectletter.
  \end{quizsolution}

  \begin{quizquestion}[
    id=ow-sound-022,
    topic=sound-intensity,
    difficulty=applied,
    marks=1,
    correct=A,
    tags={intensity,decibels},
    outcome={Convert intensity ratio to sound-level change}
  ]
    If a sound intensity increases by a factor of \(1000\), the sound level
    increases by
    \choices
      {\(\SI{30}{dB}\)}
      {\(\SI{20}{dB}\)}
      {\(\SI{10}{dB}\)}
      {\(\SI{3}{dB}\)}
      {\(\SI{1000}{dB}\)}
  \end{quizquestion}
  \begin{quizsolution}
    \(\Delta\beta=10\log_{10}(1000)=10(3)=\SI{30}{dB}\). The correct answer is
    \quizcorrectletter.
  \end{quizsolution}

  \begin{quizquestion}[
    id=ow-sound-023,
    topic=sound-intensity,
    difficulty=foundation,
    marks=1,
    correct=D,
    tags={inverse-square},
    outcome={Apply inverse-square spreading for a point sound source}
  ]
    If distance from a point sound source is doubled, the intensity becomes
    \choices
      {twice as large.}
      {half as large.}
      {unchanged.}
      {one fourth as large.}
      {four times as large.}
  \end{quizquestion}
  \begin{quizsolution}
    For a point source, \(I=P_s/(4\pi r^2)\). Doubling \(r\) makes intensity
    \(1/4\) as large. The correct answer is \quizcorrectletter.
  \end{quizsolution}

  \begin{quizquestion}[
    id=ow-pipe-024,
    topic=sound-pipes,
    difficulty=applied,
    marks=1,
    correct=B,
    tags={closed-pipe,harmonics},
    outcome={Compute closed-open pipe resonance}
  ]
    A pipe closed at one end has length \(\SI{0.85}{m}\). What is its
    fundamental frequency, using \(v=\SI{343}{m.s^{-1}}\)?
    \choices
      {\(\SI{50}{Hz}\)}
      {\(\SI{101}{Hz}\)}
      {\(\SI{202}{Hz}\)}
      {\(\SI{303}{Hz}\)}
      {\(\SI{404}{Hz}\)}
  \end{quizquestion}
  \begin{quizsolution}
    For a closed-open pipe, \(f_1=v/(4L)=343/[4(0.85)]=\SI{101}{Hz}\). The
    correct answer is \quizcorrectletter.
  \end{quizsolution}

  \begin{quizquestion}[
    id=ow-pipe-025,
    topic=sound-pipes,
    difficulty=foundation,
    marks=1,
    correct=E,
    tags={closed-pipe,odd-harmonics},
    outcome={Recall allowed harmonics in a closed-open pipe}
  ]
    Which harmonics are allowed in a pipe that is closed at one end and open at
    the other?
    \choices
      {Only even harmonics}
      {All integer harmonics}
      {Only the second harmonic}
      {Only the zeroth harmonic}
      {Only odd harmonics}
  \end{quizquestion}
  \begin{quizsolution}
    A closed-open pipe has \(n=1,3,5,\ldots\), so only odd harmonics are allowed.
    The correct answer is \quizcorrectletter.
  \end{quizsolution}

  \begin{quizquestion}[
    id=ow-beats-026,
    topic=beats,
    difficulty=applied,
    marks=1,
    correct=C,
    tags={beats},
    outcome={Compute beat frequency}
  ]
    Two tuning forks have frequencies \(\SI{440}{Hz}\) and \(\SI{446}{Hz}\). What
    beat frequency is heard?
    \choices
      {\(\SI{3}{Hz}\)}
      {\(\SI{4}{Hz}\)}
      {\(\SI{6}{Hz}\)}
      {\(\SI{443}{Hz}\)}
      {\(\SI{886}{Hz}\)}
  \end{quizquestion}
  \begin{quizsolution}
    \(f_{\text{beat}}=|446-440|=\SI{6}{Hz}\). The correct answer is
    \quizcorrectletter.
  \end{quizsolution}

  \begin{quizquestion}[
    id=ow-doppler-027,
    topic=doppler-effect,
    difficulty=applied,
    marks=1,
    correct=A,
    tags={doppler,moving-source},
    outcome={Compute observed frequency for an approaching source}
  ]
    A stationary detector hears a source moving toward it at
    \(\SI{30}{m.s^{-1}}\). The source emits \(\SI{800}{Hz}\). What frequency is
    observed?
    \choices
      {\(\SI{877}{Hz}\)}
      {\(\SI{800}{Hz}\)}
      {\(\SI{735}{Hz}\)}
      {\(\SI{343}{Hz}\)}
      {\(\SI{30}{Hz}\)}
  \end{quizquestion}
  \begin{quizsolution}
    For an approaching source and stationary detector,
    \(f'=f v/(v-v_S)=800(343)/(343-30)=\SI{877}{Hz}\). The correct answer is
    \quizcorrectletter.
  \end{quizsolution}

  \begin{quizquestion}[
    id=ow-doppler-028,
    topic=doppler-effect,
    difficulty=foundation,
    marks=1,
    correct=D,
    tags={doppler,concept},
    outcome={Predict qualitative Doppler shift}
  ]
    If a sound source moves away from a stationary listener, the observed
    frequency is
    \choices
      {higher than the emitted frequency.}
      {always doubled.}
      {equal to the emitted frequency.}
      {lower than the emitted frequency.}
      {zero.}
  \end{quizquestion}
  \begin{quizsolution}
    Motion away stretches the wavefront spacing at the listener, so the observed
    frequency is lower. The correct answer is \quizcorrectletter.
  \end{quizsolution}

  \begin{quizquestion}[
    id=ow-shock-029,
    topic=shock-waves,
    difficulty=foundation,
    marks=1,
    correct=B,
    tags={mach-cone},
    outcome={State the condition for a shock wave}
  ]
    A shock wave forms when a source
    \choices
      {moves slower than the wave speed.}
      {moves faster than the sound speed in the medium.}
      {has zero frequency.}
      {emits only transverse waves.}
      {is stationary.}
  \end{quizquestion}
  \begin{quizsolution}
    Shock waves occur when the source speed exceeds the sound speed in the
    medium. The correct answer is \quizcorrectletter.
  \end{quizsolution}

  \begin{quizquestion}[
    id=ow-summary-030,
    topic=summary,
    difficulty=foundation,
    marks=1,
    correct=E,
    tags={strategy,units},
    outcome={Choose a reliable problem-solving first step}
  ]
    What is the best first step when a problem mixes oscillations and waves?
    \choices
      {Assume every frequency is angular frequency.}
      {Use the Doppler formula first.}
      {Ignore units until the final answer.}
      {Treat all waves as standing waves.}
      {Identify whether the unknown belongs to one oscillator, a traveling wave, a standing wave, or sound.}
  \end{quizquestion}
  \begin{quizsolution}
    The formula family depends on the physical situation. First classify the
    situation, then choose the equation. The correct answer is
    \quizcorrectletter.
  \end{quizsolution}
\end{quizbank}

\quizdefineversion{A}{%
  \quizselectall
  \quizshuffleoptions{151617}%
}

\quizuseversion{A}

\pagenumbering{gobble}
\makequiztitle
\pagenumbering{arabic}
\tableofcontents
\clearpage

\printquizquestions[2]

\begin{quizanswerkeycontent}
  \clearpage
  \printquizanswerkey
\end{quizanswerkeycontent}

\begin{quizsolutioncontent}
  \clearpage
  \printquizsolutions
\end{quizsolutioncontent}

\begin{quizteachercontent}
  \clearpage
  \printquiztopicreport
  \printquizteacherreport
\end{quizteachercontent}

\end{document}
